% New presentation draft: storyline and content first, polish later by switching to \usetheme[progress]{presentation}. Build: latexmk -pdf
\documentclass[aspectratio=169,11pt]{beamer}
\usetheme{draft}

\title{Deck title}
\subtitle{What this talk gives you}
\author{Author}
\institute{Team or institution}% optional
\date{Venue \textperiodcentered{} Year}

\begin{document}

\begin{frame}[plain]
  \titlepage
\end{frame}

\begin{frame}{Agenda}{Question of the talk}
  \tableofcontents
  \note{The question the talk answers and why the chapters come in this order.}
\end{frame}

% Every \section opens a chapter slide, its subsections become that chapter's agenda
\section{First chapter}
\subsection{First step}

\begin{frame}{Claim as a full sentence}
  \begin{itemize}
    \item Evidence or argument for the claim
    \item What follows from it
  \end{itemize}
  \visual{Chart: what it shows and what it proves}
  \question{What is still open on this slide}
  \note{Why this claim matters and what the audience should take from it.}
\end{frame}

\section{Second chapter}
\subsection{Second step}

\begin{frame}{Claim as a full sentence}
  \begin{itemize}
    \item Evidence or argument for the claim
  \end{itemize}
  \note{How this chapter builds on the first one.}
\end{frame}

\begin{frame}[plain,c]
  \centering\Large The one sentence to remember.
  \note{The one sentence again in your own words, then the next step for the audience.}
\end{frame}

\end{document}
